% ============================================================
% Detailed draft for Shanjun / Thread B
% Scope: Gemma-2 experiments only
% Paste this block before \bibliographystyle{plainnat} and \bibliography{refs}.
% ============================================================

\section*{Draft: Shanjun's Gemma-2 Experiments}

\paragraph{Role in the joint paper.}
This section summarizes the Gemma-2 experiments conducted in the SDForger-style
HAR time-series generation pipeline. In the joint paper plan, these experiments
mainly support RQ2: \emph{which conditioning knobs help, and what trade-offs do
they introduce?} The experiments also provide supporting evidence for RQ1 by
showing that syntactically valid generations can still fail in class balance,
similarity, or downstream utility, and they motivate RQ3 by showing that channel
selection, statistical prompting, latent dimensionality, and filtering do not by
themselves fully remove collapse.

\subsection*{Chronological Development of the Gemma-2 Experiment Line}

\paragraph{Project branches and experiment stages.}
The Git history in the SDForger-PAMAP2 repository mirrors the experiment
progression. The early archive branches store the two-activity Gemma results,
the MM-Fit branch adds the second-dataset pipeline, and the current
stats/filtering branches synchronize the window-statistics and filtering
results. Table~\ref{tab:shanjun-branch-timeline} lists the stages that are
relevant to this draft.

\begin{table}[t]
\centering
\small
\caption{Chronological organization of Shanjun's Gemma-2 experiment line.}
\label{tab:shanjun-branch-timeline}
\resizebox{\linewidth}{!}{%
\begin{tabular}{llll}
\toprule
Stage & Main question & Main files / branch evidence & Report role \\
\midrule
Early PAMAP2 Gemma & Does the SDForger-style pipeline run with Gemma-2? &
\texttt{gemma}, \texttt{gemma\_5class}, \texttt{gemma\_5class\_v2} outputs &
Setup / weak-conditioning evidence \\
Two-activity PAMAP2 & Can walking/running simplify the analysis? &
\texttt{codex/archive-pamap2-experiments} &
Main RQ2 setting \\
Channel sweep & Which sensor subset exposes useful structure? &
\texttt{gemma\_2act\_\{full,hand,chest,...\}} &
Main RQ2 result \\
Stats prompt and epoch sweep & Do window statistics and training length help? &
\texttt{codex/archive-two-act-stats-results} &
Main RQ2 result and negative ablation \\
ICA sweep & Does latent dimensionality improve controllability? &
\texttt{gemma\_2act\_*stats\_ica\{3,6,12,18,24\}} &
RQ2 ablation / bridge to RQ3 \\
Window-128 and filtering & Do granularity and post-hoc QC help? &
\texttt{codex/sync-256-stats-filter-*} &
Secondary ablation / appendix \\
MM-Fit extension & Does the same pipeline transfer to a second dataset? &
\texttt{codex/mmfit-pipeline} &
External-validity / appendix \\
Raw-window Gemma probe & Can we bypass FastICA with raw sequence serialization? &
\texttt{pamap2\_gemma2\_2b\_2act\_chest\_gyro\_raw\_*} &
Diagnostic / negative evidence \\
\bottomrule
\end{tabular}}
\end{table}

\subsection*{Datasets and Preprocessing Choices}

\paragraph{PAMAP2.}
The main experiments use PAMAP2 with subject-level splits. The five-class
setting uses \texttt{walking}, \texttt{running}, \texttt{cycling},
\texttt{sitting}, and \texttt{standing}. The main two-activity setting keeps only
\texttt{walking} and \texttt{running}, because this reduces label ambiguity and
focuses the analysis on periodic gait-like motion. Unless otherwise stated, the
two-activity experiments use train subjects 101--106, validation subject 107,
and test subject 108. The standard window size is 256 with stride 128; a later
secondary experiment uses window size 128 and stride 64.

\paragraph{Sensor channel subsets.}
The experiments deliberately vary the selected PAMAP2 channels. The full
setting uses 18 channels from hand, chest, and ankle accelerometer/gyroscope
signals. The two-activity Gemma experiments then compare smaller channel groups:
\texttt{hand}, \texttt{chest}, \texttt{hand\_acc}, \texttt{chest\_acc},
\texttt{hand\_gyro}, and \texttt{chest\_gyro}. This channel choice is treated as
a representation/conditioning knob rather than a mere preprocessing detail,
because it changes which structure is exposed to the compact latent code.

\paragraph{Latent representation.}
Most experiments use FastICA as the compact latent representation. The baseline
two-activity runs use 6 ICA components per selected channel. The ICA sweep tests
3, 6, 12, 18, and 24 components in selected channel/statistics settings. A
later diagnostic line uses \texttt{raw\_window} serialization for
\texttt{chest\_gyro} to test whether bypassing FastICA avoids representation
loss; this is not used as a main result because it currently lacks complete
similarity/utility evaluation and shows class-coverage issues.

\paragraph{MM-Fit.}
The MM-Fit extension uses five exercise classes: \texttt{squats},
\texttt{pushups}, \texttt{lunges}, \texttt{bicep\_curls}, and
\texttt{jumping\_jacks}. The selected streams are left/right smartwatch
accelerometer and gyroscope signals, yielding 12 channels. Windows are 5 seconds
long, sampled to 250 time steps, with 2.5-second stride. The split follows
workout-session groups rather than PAMAP2 subject IDs. MM-Fit is therefore used
as a second-dataset sanity check, not for direct numerical comparison with
PAMAP2.

\begin{table}[t]
\centering
\small
\caption{Dataset processing choices for the Gemma-2 experiment line.}
\label{tab:shanjun-datasets}
\resizebox{\linewidth}{!}{%
\begin{tabular}{llllll}
\toprule
Dataset / setting & Activities & Channels & Window / stride & Representation & Main role \\
\midrule
PAMAP2 5-class & walk/run/cycle/sit/stand & 18 or selected subsets & 256 / 128 & FastICA & Early baseline \\
PAMAP2 2-act & walking/running & full, hand, chest, acc/gyro subsets & 256 / 128 & FastICA & Main RQ2 setting \\
PAMAP2 2-act win128 & walking/running & hand, hand\_acc, chest\_gyro & 128 / 64 & FastICA & Secondary ablation \\
PAMAP2 raw-window & walking/running & chest\_gyro only & 256 / 128 & raw sequence tokens & Diagnostic probe \\
MM-Fit 5-class & squat/pushup/lunge/curl/jumping jack & L/R watch acc+gyro & 250 / 125 & FastICA & Second dataset \\
\bottomrule
\end{tabular}}
\end{table}

\subsection*{Gemma-2 Model and Shared Evaluation Protocol}

All main Gemma experiments use \texttt{google/gemma-2-2b} with LoRA fine-tuning.
The typical LoRA configuration uses rank 16, alpha 32, dropout 0.05, and targets
attention and MLP projection modules. The evaluation reports three types of
outputs:
\begin{itemize}
  \item generation validity: prompt count, parse success, valid generated count,
  duplicate ratio, and generated class counts;
  \item similarity/fidelity: MDD, ACD, SD, KD, ED, DTW, and SHR, where SD is
  Skewness Difference in the SDForger/TSGBench sense rather than Spectral
  Difference; ED/real and DTW/real ratios are computed against a real-vs-real
  baseline when available;
  \item downstream HAR utility: real-only, synthetic-only, and real-plus-synthetic
  RandomForest accuracy and macro-F1.
\end{itemize}
For the final report, all numerical comparisons should be interpreted within the
same protocol. ED and DTW should not be compared as absolute numbers across
different channel normalizations or datasets.

\paragraph{Table-display rule after the metric update.}
Because the complete similarity block now contains seven metrics
(MDD/ACD/SD/KD/ED/DTW/SHR), the main-text Gemma tables should remain compact
and should not repeat every metric in every ablation table. In the main body, we
keep validity, ED/real, DTW/real, and HAR utility as the most readable summary
columns. We then add one full similarity table, either in the main text for the
representative settings or in the appendix for the full sweep, with columns
\texttt{MDD}, \texttt{ACD}, \texttt{SD (skew.)}, \texttt{KD}, \texttt{ED},
\texttt{DTW}, and \texttt{SHR}. This keeps the report aligned with
SDForger-style evaluation while avoiding very wide main-result tables.

\subsection*{Stage 1: Early PAMAP2 Five-Class Gemma Runs}

The first Gemma experiments tested whether the SDForger-style pipeline could be
trained and evaluated beyond GPT-2. These runs used the five-class PAMAP2 setup
and later variants with selected channel subsets. They are useful background,
but they should not dominate the main results because the protocols changed
across runs.

\begin{table}[t]
\centering
\small
\caption{Selected early PAMAP2 five-class Gemma results. These are supporting
evidence, not the main RQ2 result. ``Cond. acc.'' refers to the extended
condition-consistency classifier, not the TSTR utility classifier.}
\label{tab:shanjun-fiveclass}
\resizebox{\linewidth}{!}{%
\begin{tabular}{lrrrrrr}
\toprule
Run & Valid synthetic & W-dist. gap & DTW gap & Cond. acc. & Cond. macro-F1 & Suggested use \\
\midrule
\texttt{gemma} & 227 & 1.303 & 296.631 & 0.815 & 0.261 & early baseline \\
\texttt{gemma\_5class} & 636 & 1.926 & 320.165 & 0.310 & 0.166 & weak conditioning \\
\texttt{gemma\_5class\_v2} & 857 & 1.629 & 258.321 & 0.622 & 0.581 & improved baseline \\
\texttt{gemma\_5class\_v2\_chest} & 695 & 0.994 & 96.133 & 0.230 & 0.136 & fidelity-oriented subset \\
\texttt{gemma\_5class\_v2\_hand} & 362 & 2.583 & 486.448 & 0.312 & 0.237 & weak utility \\
\texttt{gemma\_5class\_v2\_hand\_chest} & 857 & 1.506 & 271.590 & 0.685 & 0.666 & best 5-class condition consistency \\
\bottomrule
\end{tabular}}
\end{table}

\paragraph{Brief interpretation.}
These runs show that switching to Gemma-2 and changing channel subsets can
produce syntactically valid synthetic windows, but the relationship between
fidelity and condition consistency is unstable. This motivated the later
two-activity setting, where the channel/prompt trade-off could be studied more
cleanly.

\subsection*{Stage 2: PAMAP2 Two-Activity Channel Sweep}

The main experiment line restricts PAMAP2 to \texttt{walking} and
\texttt{running}. The first two-activity sweep compares sensor channel subsets
without statistical prompts. The goal is to test whether channel selection
itself is a conditioning/representation knob.

\begin{table}[t]
\centering
\small
\caption{PAMAP2 two-activity Gemma channel sweep without statistical prompts.
ED/real and DTW/real are computed against the corresponding real-vs-real
baseline for each channel group.}
\label{tab:shanjun-channel}
\resizebox{\linewidth}{!}{%
\begin{tabular}{lrrrrrrrr}
\toprule
Setting & Valid & Run & Walk & ED/real & DTW/real & Syn acc. & Syn F1 & Real+Syn acc. \\
\midrule
\texttt{full} & 83/305 & 27 & 56 & 3.084 & 5.678 & 0.390 & 0.331 & 0.984 \\
\texttt{hand} & 226/305 & 82 & 144 & 4.158 & 10.162 & 0.921 & 0.917 & 0.977 \\
\texttt{chest} & 213/305 & 102 & 111 & 0.862 & 2.206 & 0.384 & 0.297 & 0.964 \\
\texttt{hand\_acc} & 181/305 & 71 & 110 & 4.275 & 11.272 & 0.387 & 0.303 & 0.961 \\
\texttt{chest\_acc} & 274/305 & 110 & 164 & 0.826 & 2.410 & 0.220 & 0.204 & 0.964 \\
\texttt{hand\_gyro} & 202/305 & 94 & 108 & 1.104 & 3.902 & 0.387 & 0.303 & 0.689 \\
\texttt{chest\_gyro} & 249/305 & 101 & 148 & 0.721 & 1.130 & 0.374 & 0.291 & 0.954 \\
\bottomrule
\end{tabular}}
\end{table}

\paragraph{Brief interpretation.}
The hand-channel setting is strongest for downstream synthetic-only utility,
whereas chest gyroscope is strongest for similarity. Full-channel generation is
not automatically best. This supports the joint paper's claim that what the
latent representation exposes matters as much as the nominal LLM capacity.

\subsection*{Stage 3: Epoch Scaling Without Statistical Prompts}

We next tested whether longer training alone improves generation. Selected
strong or informative channel groups were trained for 5, 10, and 20 epochs.

\begin{table}[t]
\centering
\small
\caption{Epoch scaling without statistical prompts. Longer training changes the
trade-off but does not monotonically improve utility or similarity.}
\label{tab:shanjun-epoch}
\resizebox{\linewidth}{!}{%
\begin{tabular}{lrrrrrrr}
\toprule
Setting & Epoch & Valid & ED/real & DTW/real & Syn acc. & Syn F1 & Real+Syn acc. \\
\midrule
\texttt{hand} & 5 & 226/305 & 4.158 & 10.162 & 0.921 & 0.917 & 0.977 \\
\texttt{hand} & 10 & 75/305 & 4.584 & 11.115 & 0.820 & 0.817 & 0.977 \\
\texttt{hand} & 20 & 139/305 & 4.378 & 10.746 & 0.377 & 0.297 & 0.974 \\
\texttt{chest\_gyro} & 5 & 249/305 & 0.721 & 1.130 & 0.374 & 0.291 & 0.954 \\
\texttt{chest\_gyro} & 10 & 294/305 & 0.454 & 0.736 & 0.370 & 0.286 & 0.954 \\
\texttt{chest\_gyro} & 20 & 301/305 & 0.653 & 1.065 & 0.370 & 0.286 & 0.961 \\
\texttt{chest} & 5 & 213/305 & 0.862 & 2.206 & 0.384 & 0.297 & 0.964 \\
\texttt{chest} & 10 & 298/305 & 0.944 & 1.684 & 0.380 & 0.295 & 0.970 \\
\texttt{chest} & 20 & 74/305 & 1.030 & 2.904 & 0.626 & 0.626 & 0.964 \\
\bottomrule
\end{tabular}}
\end{table}

\paragraph{Brief interpretation.}
More epochs are not a reliable fix. For example, \texttt{hand} remains useful at
5 epochs but degrades at 20 epochs, while \texttt{chest\_gyro} improves
similarity at 10 epochs without improving synthetic-only utility. This is a
negative ablation: training longer changes the model behavior but does not
resolve collapse.

\subsection*{Stage 4: Statistical Prompt Conditioning}

The next intervention adds window-level statistics to the serialized training
text: mean, standard deviation, minimum, and maximum, with three decimal digits.
The purpose is to strengthen the conditioning signal by explicitly exposing
summary statistics to the LLM.

\begin{table}[t]
\centering
\small
\caption{Statistical prompt results for selected channel groups and epochs.}
\label{tab:shanjun-stats}
\resizebox{\linewidth}{!}{%
\begin{tabular}{lrrrrrrr}
\toprule
Setting & Epoch & Valid & ED/real & DTW/real & Syn acc. & Syn F1 & Real+Syn acc. \\
\midrule
\texttt{hand + stats} & 5 & 301/305 & 3.912 & 9.288 & 0.928 & 0.925 & 0.970 \\
\texttt{hand + stats} & 10 & 261/305 & 4.561 & 10.379 & 0.846 & 0.844 & 0.970 \\
\texttt{hand + stats} & 20 & 284/305 & 4.445 & 10.349 & 0.921 & 0.919 & 0.957 \\
\texttt{chest\_gyro + stats} & 5 & 283/305 & 0.466 & 0.740 & 0.449 & 0.412 & 0.954 \\
\texttt{chest\_gyro + stats} & 10 & 303/305 & 0.522 & 0.795 & 0.374 & 0.298 & 0.957 \\
\texttt{chest\_gyro + stats} & 20 & 271/305 & 0.622 & 1.016 & 0.564 & 0.554 & 0.957 \\
\texttt{chest + stats} & 5 & 270/305 & 0.692 & 1.556 & 0.370 & 0.293 & 0.967 \\
\texttt{chest\_acc + stats} & 5 & 280/305 & 0.858 & 2.307 & 0.370 & 0.296 & 0.967 \\
\bottomrule
\end{tabular}}
\end{table}

\paragraph{Brief interpretation.}
Statistical prompts generally improve validity and can improve selected utility
or similarity, but they do not remove the trade-off. \texttt{hand + stats} is the
strongest utility-oriented setting, while \texttt{chest\_gyro + stats} is the
strongest similarity-oriented setting. This suggests that prompt statistics help
only when the selected channels and latent representation expose information
that the model can steer.

\subsection*{Stage 5: ICA Dimensionality Sweep}

To test whether latent capacity is the bottleneck, we varied the number of
FastICA components per channel in the two strongest statistics-prompt families:
\texttt{chest\_gyro + stats} and \texttt{hand + stats}.

\begin{table}[t]
\centering
\small
\caption{Condensed ICA sweep. The table keeps representative rows that should be
shown in the main report; the full sweep can be placed in an appendix.}
\label{tab:shanjun-ica}
\resizebox{\linewidth}{!}{%
\begin{tabular}{llrrrrrr}
\toprule
Setting & ICA / epoch & Valid & ED/real & DTW/real & Syn acc. & Syn F1 & Real+Syn acc. \\
\midrule
\texttt{chest\_gyro + stats} & ICA3 / ep5 & 253/305 & 0.734 & 1.206 & 0.377 & 0.290 & 0.957 \\
\texttt{chest\_gyro + stats} & ICA6 / ep5 & 283/305 & 0.466 & 0.740 & 0.449 & 0.412 & 0.954 \\
\texttt{chest\_gyro + stats} & ICA12 / ep5 & 302/305 & 0.569 & 0.825 & 0.898 & 0.894 & 0.964 \\
\texttt{chest\_gyro + stats} & ICA12 / ep10 & 299/305 & 0.758 & 0.831 & 0.964 & 0.961 & 0.964 \\
\texttt{chest\_gyro + stats} & ICA18 / ep10 & 236/305 & 0.624 & 1.084 & 0.859 & 0.841 & 0.951 \\
\texttt{hand + stats} & ICA6 / ep5 & 301/305 & 3.912 & 9.288 & 0.928 & 0.925 & 0.970 \\
\texttt{hand + stats} & ICA12 / ep10 & 258/305 & 4.058 & 9.929 & 0.928 & 0.925 & 0.974 \\
\texttt{hand + stats} & ICA18 / ep10 & 294/305 & 3.908 & 9.015 & 0.954 & 0.952 & 0.974 \\
\bottomrule
\end{tabular}}
\end{table}

\paragraph{Brief interpretation.}
Increasing ICA dimensionality can help, but the effect is not monotonic.
\texttt{chest\_gyro + stats} benefits strongly from ICA12 in downstream utility,
whereas similarity is already strong at ICA6. This supports the distinction
between representation capacity and steerability: adding dimensions helps only
when the new latent variables align with conditionable structure.

\subsection*{Stage 6: Window-128, Statistics Source, and Filtering}

The window-128 experiments increase the number of training/test windows and test
whether shorter windows and alternative statistics sources improve generation.
These runs are valuable but should be treated as secondary ablations because the
protocol differs from the main window-256 setting.

\begin{table}[t]
\centering
\small
\caption{Window-128 statistics-prompt experiments. These results are useful as a
secondary ablation because the window length and number of prompts differ from
the main setting. ED/ref and DTW/ref are retained as compact distance summaries;
the full SDForger-style similarity metrics, including SD (skew.) and SHR, should
be reported in the accompanying full-metric CSV or appendix table.}
\label{tab:shanjun-win128}
\resizebox{\linewidth}{!}{%
\begin{tabular}{lrrrrrrr}
\toprule
Setting & Source & Valid rate & ED/ref & DTW/ref & Syn acc. & Real acc. & Real+Syn acc. \\
\midrule
\texttt{hand stats ICA12 ep5} & raw stats & 0.461 & 15.613 & 24.999 & 0.952 & 0.977 & 0.971 \\
\texttt{hand stats ICA12 ep10} & raw stats & 0.032 & 15.883 & 24.906 & 0.947 & 0.977 & 0.974 \\
\texttt{hand\_acc stats ICA12 ep5} & raw stats & 0.545 & 15.074 & 30.364 & 0.926 & 0.976 & 0.955 \\
\texttt{hand\_acc stats ICA12 ep10} & raw stats & 0.976 & 15.143 & 31.652 & 0.745 & 0.976 & 0.961 \\
\texttt{chest\_gyro stats ICA12 ep5} & raw stats & 0.992 & 1.314 & 1.809 & 0.503 & 0.961 & 0.958 \\
\texttt{chest\_gyro stats ICA12 ep10} & raw stats & 0.850 & 0.891 & 1.312 & 0.357 & 0.961 & 0.955 \\
\texttt{chest\_gyro icastats ep5} & embedding stats & 0.980 & -- & -- & 0.955 & 0.961 & 0.964 \\
\texttt{chest\_gyro icastats ep10} & embedding stats & 0.919 & -- & -- & 0.959 & 0.961 & 0.962 \\
\bottomrule
\end{tabular}}
\end{table}

\paragraph{Post-hoc filtering.}
The p25/p50/p75 consistency-filtering experiments keep only a target fraction of
generated candidates according to a statistics-consistency score. They are useful
for quality control, but the results should be placed in an appendix: filtering
can improve selected downstream numbers or maintain real-plus-synthetic accuracy,
but it reduces the number of retained synthetic samples and does not solve the
underlying generation problem. After the SDForger-style recomputation, the
filtering table should additionally report the full similarity block in the
appendix, because filtering can change distributional metrics such as MDD, SD
(skew.), KD, and SHR even when the compact ED/ref and DTW/ref columns look
similar.

\paragraph{Brief interpretation.}
Window granularity and statistics source change the validity/utility/similarity
trade-off. In particular, embedding-channel statistics for \texttt{chest\_gyro}
give high synthetic-only utility in the window-128 setting, but the protocol is
different enough that it should be framed as a secondary finding rather than as
the main result.

\subsection*{Stage 7: Raw-Window Gemma Diagnostic}

A later diagnostic experiment replaces FastICA latents with serialized raw
\texttt{chest\_gyro} windows. The motivation is to test whether compact linear
latent codes are the main bottleneck. The raw-window representation uses 768
values per window (256 time steps times 3 channels), per-channel standardization,
and window-level statistics in the prompt.

\begin{table}[t]
\centering
\small
\caption{Raw-window Gemma diagnostic probes. These runs are not main results
because they currently lack complete similarity/utility evaluation and show
class-coverage issues.}
\label{tab:shanjun-raw}
\resizebox{\linewidth}{!}{%
\begin{tabular}{lrrrrll}
\toprule
Run & Prompts & Parse success & Valid & Token budget & Generated classes & Suggested use \\
\midrule
\texttt{genprobe8} & 8 & 8 & 0 & 2300 & none & negative diagnostic \\
\texttt{genprobe8\_tok6500} & 8 & 8 & 8 & 6500 & walking only & parse/length diagnostic \\
\texttt{genprobe32\_tok6500} & 32 & 32 & 32 & 6500 & walking only & class-coverage diagnostic \\
\bottomrule
\end{tabular}}
\end{table}

\paragraph{Brief interpretation.}
Longer token budgets make raw-window parsing possible, but the valid candidates
cover only \texttt{walking} in the observed probes. This suggests that simply
removing FastICA and serializing raw windows is not yet a reliable solution; it
also makes the generation problem much longer and harder for the LLM.

\subsection*{Stage 8: MM-Fit Gemma Extension}

The MM-Fit experiment tests whether the pipeline can be transferred to another
HAR-like dataset. It uses 12 smartwatch channels and five exercise classes. This
experiment should be reported as a second-dataset sanity check rather than a
direct comparison against PAMAP2.

\begin{table}[t]
\centering
\small
\caption{MM-Fit Gemma-2 five-class result. This is a second-dataset extension,
not directly comparable to PAMAP2 absolute values.}
\label{tab:shanjun-mmfit}
\resizebox{\linewidth}{!}{%
\begin{tabular}{lrrrrrrrr}
\toprule
Dataset & Prompts & Valid & Classes covered & ED & DTW & Syn acc. & Syn F1 & Real+Syn acc. \\
\midrule
MM-Fit Gemma-2 & 406 & 369 & 5/5 & 86.460 & 218.915 & 0.515 & 0.538 & 0.975 \\
\bottomrule
\end{tabular}}
\end{table}

\paragraph{Brief interpretation.}
The MM-Fit Gemma run demonstrates that the pipeline can be moved beyond PAMAP2,
with full class coverage and non-empty utility evaluation. However, the
real-only classifier is already very strong on MM-Fit, and the synthetic-only
score remains far below real-only. This supports the broader thesis that LLM
generation requires stronger conditioning and representation design, not only
dataset transfer.

\subsection*{Recommended Results to Show in the Final Report}

For the main body, I recommend showing the following results:
\begin{itemize}
  \item the two-activity PAMAP2 channel sweep
  (Table~\ref{tab:shanjun-channel});
  \item the epoch-scaling negative ablation
  (Table~\ref{tab:shanjun-epoch});
  \item the statistical-prompt comparison
  (Table~\ref{tab:shanjun-stats});
  \item a condensed ICA-dimensionality table
  (Table~\ref{tab:shanjun-ica});
  \item one secondary table for window-128/statistics-source experiments if space
  allows (Table~\ref{tab:shanjun-win128}).
\end{itemize}
The early five-class Gemma table, MM-Fit extension, raw-window diagnostic, and
filtering results should be placed in the appendix or summarized briefly in the
main text. In addition, the final report should include one appendix table or
machine-readable supplementary CSV for the complete SDForger-style similarity
metrics, because the compact main tables intentionally omit MDD, ACD, SD
(skew.), KD, and SHR for readability.

\subsection*{Short Result Analysis}

Across the Gemma-2 experiments, the most consistent observation is that no single
knob dominates all metrics. Channel choice controls what the latent code exposes:
\texttt{hand} produces stronger downstream utility, while \texttt{chest\_gyro}
produces stronger similarity. Statistical prompts strengthen the conditioning
interface and improve validity, but their gains depend on the selected channel
and latent dimensionality. Longer training is not a reliable fix; it can improve
some similarity metrics while leaving synthetic-only utility unchanged or worse.
Increasing ICA dimensionality helps locally, especially for
\texttt{chest\_gyro + stats}, but the effect is non-monotonic. Post-hoc
filtering and raw-window serialization are useful diagnostics, yet neither
provides a complete solution. These findings support the joint paper's central
claim: controllability in LLM-based synthetic HAR generation is primarily a
conditioning-and-representation problem, not simply a matter of using a larger
model or training longer.
